\section{Canonical classes and magnitude homology}
\label{sec:canonical}

Throughout the paper $M$ is a connected closed negatively curved
Riemannian manifold of dimension $n\geq2$, $\Gamma=\pi_1(M)$, and
$G=\Gamma\times\Gamma$. We use integral coefficients unless indicated
otherwise. The definitions in this section require considerably less
geometry than negative curvature.

\subsection{The Costa--Farber class}

The endpoint fibration of an aspherical space can be compared with
the covering corresponding to the diagonal subgroup
$\Delta\Gamma\leq G$; see \cite{CostaFarber2010}.
Let $P=\Z[\Gamma]$ with the $G$-action
\[
 (g,h)a=gah^{-1}.
\]
Write $I=\ker(\varepsilon:P\to\Z)$.
The augmentation sequence is an extension of $G$-modules:
\begin{equation}\label{eq:augmentation}
 0\longrightarrow I\longrightarrow P
 \xrightarrow{\varepsilon}\Z\longrightarrow0.
\end{equation}
Its connecting class is
$\mathfrak v\in H^1(G;I)$. With the inhomogeneous convention
for group cochains, $\mathfrak v$ is represented by
\begin{equation}\label{eq:canonical-cocycle}
 c(g,h)=(g,h)\cdot1-1=gh^{-1}-1.
\end{equation}
In particular, the restriction of $\mathfrak v$ to
$\Delta\Gamma$ vanishes.

For each $k\geq1$, let
$q_k:I^{\otimes k}\to\bigwedge^k I$ be the quotient
mapping pure tensors to exterior products. Define
\begin{equation}\label{eq:exterior-class}
 \mathfrak a_k=q_{k*}(\mathfrak v^k)
 \in H^k(G;\bigwedge^k I).
\end{equation}
We emphasize that $q_k$ is the exterior quotient, not the
unnormalized alternating sum. No factorial occurs in
\eqref{eq:exterior-class}.

\subsection{The full simplex and exterior augmentation}

Let $Y$ be the full simplicial complex on the vertex set $\Gamma$.
Its oriented simplicial chain group in degree $k$ is
\[
 C_k(Y;\Z)=\bigwedge^{k+1}P.
\]
The augmentation of $P$ induces a contraction
\[
 \iota_\varepsilon:\bigwedge^{k+1}P
 \longrightarrow\bigwedge^kP,\quad
 \iota_\varepsilon(a_0\wedge\cdots\wedge a_k)
 =\sum_{i=0}^{k}(-1)^i
   a_0\wedge\cdots\widehat{a_i}\cdots\wedge a_k,
\]
which is exactly the simplicial boundary.

\begin{lemma}\label{lem:exterior-quotient}
For $k\geq1$ there is a natural $G$-equivariant exact sequence
\[
 C_{k+1}(Y)\xrightarrow{\partial}C_k(Y)
 \xrightarrow{T_k}\bigwedge^kI\longrightarrow0,
\]
where
\[
 T_k[a_0,\ldots,a_k]
 =(a_1-a_0)\wedge\cdots\wedge(a_k-a_0).
\]
Consequently
$C_k(Y)/\partial C_{k+1}(Y)\cong\bigwedge^kI$.
\end{lemma}

\begin{proof}
Choose any $a_*\in\Gamma$ and split the underlying abelian
group as $P=\Z a_*\oplus I$. Then
\[
 \bigwedge^{k+1}P
 = (a_*\wedge\bigwedge^kI)\oplus\bigwedge^{k+1}I.
\]
Contraction with $\varepsilon$ sends $a_*\wedge w$ to $w$
and vanishes on $\bigwedge^{k+1}I$.
The image of the contraction in degree $k+1$ is
$\bigwedge^{k+1}I$, which is precisely the kernel of
contraction in degree $k$. Thus contraction induces the
claimed quotient. The displayed expression for $T_k$
follows by expressing all vertices as
$a_0+(a_i-a_0)$ and expanding the wedge.
Although the splitting is not equivariant, the resulting map is.
\end{proof}

Write $Q=|Y|/G$ and
$\mathcal V_k=(\bigwedge^k I)_G$.
The action has no inversions on simplices (using the unique-root
property of torsion-free hyperbolic groups
\cite{BridsonHaefliger1999}): if an element
permuted a finite set of vertices, some positive power would
fix each of them; the unique-root property for torsion-free
hyperbolic groups then forces each vertex to have been fixed
already. The cellular chains of $Q$ therefore identify with
$C_k(Y)_G$. Passing to coinvariants in the lemma gives a
map from cellular cycles to $\mathcal V_k$.

\begin{lemma}\label{lem:cycle-injective}
The assignment
\[
 \iota_k:H_k(Q;\Z)\longrightarrow\mathcal V_k,\qquad
 [z]\longmapsto[T_k(z)]
\]
is well-defined and injective for $k\geq1$.
\end{lemma}

\begin{proof}
The definition is independent of boundaries by
$T_k\partial=0$. Exactness of coinvariants applied to
Lemma~\ref{lem:exterior-quotient} identifies
$\mathcal V_k$ with
$C_k(Y)_G/\partial C_{k+1}(Y)_G$.
The homology group $H_k(Q)$ is the subgroup of that
quotient represented by cycles, proving injectivity.
\end{proof}

\begin{proposition}\label{prop:comparison}
Let $B=M\times M$ and suppose $M$ is oriented.
If $f:B\to Q$ is the map obtained from a $G$-equivariant
map $\widetilde M^2\to Y$, then, for each $k$,
the image under $\iota_k$ of the homological comparison
$f_*:H_k(B)\to H_k(Q)$ is the cap-product evaluation
of $\mathfrak a_k$, with the usual local-coefficient
identifications. In particular,
\begin{equation}\label{eq:A-comparison}
 \iota_{2n}f_*[B]
 =\mathfrak a_{2n}\frown[B].
\end{equation}
\end{proposition}

\begin{proof}
The homogeneous bar construction gives an explicit
comparison map from a free $G$-resolution to the augmented
simplex chains: a bar simplex $(g_0,\ldots,g_k)$ is sent
to $[g_0\cdot1,\ldots,g_k\cdot1]$.
Its image under $T_k$ is
\[
 (g_1\cdot1-g_0\cdot1)\wedge\cdots\wedge
 (g_k\cdot1-g_0\cdot1).
\]
The cup product of the cocycle \eqref{eq:canonical-cocycle}
evaluates on the corresponding bar simplex as the wedge
of consecutive differences. Their transformation into
the differences from the first vertex is unit triangular,
so their wedge products coincide. This proves equality
on bar chains; any equivariant map from a free resolution
is equivariantly chain-homotopic to the comparison map.
The cap-product assertion follows.
\end{proof}

\subsection{Magnitude homology at the shortest length}

We use the magnitude chain complex of Leinster--Shulman
\cite{LeinsterShulman2021}. For endpoints $x,y$ and a
length $\ell$, write $C_{k,\ell}(M;x,y)$ for the free
abelian group generated by proper chains
\[
 (x=x_0,x_1,\ldots,x_k=y),
 \qquad \sum_{i=1}^k d(x_{i-1},x_i)=\ell.
\]
The differential deletes an internal point only when it
lies metrically between its neighbors. See also the
interval-poset description in \cite{KanetaYoshinaga2021}.

\begin{proposition}\label{prop:MH-minimal}
Let $x\ne y$ in a complete negatively curved manifold $M$,
and assume the number $m_g(x,y)$ of minimizing geodesics
is finite. Then there is an isomorphism
\[
 \MH_{2,d(x,y)}(M;x,y)
 \cong\ker\bigl(\Z[S_g(u,v)]
               \xrightarrow{\varepsilon}\Z\bigr).
\]
In particular,
\[
 \operatorname{rank}\MH_{2,d(x,y)}(M;x,y)
 =m_g(x,y)-1.
\]
\end{proposition}

\begin{proof}
Write $d=d(x,y)>0$. A chain of total length $d$
must be ordered along a minimizing geodesic, by equality
in the triangle inequality. A point $z\ne x,y$ occurring
in $(x,z,y)$ determines exactly one minimizing geodesic:
the two minimizing segments splice without a corner,
because their concatenation is globally minimizing,
and the geodesic equation has unique solutions.
Conversely every minimizing geodesic contains such points.

Now $C_{1,d}(M;x,y)=\Z(x,y)$, and
\[
 \partial_2(x,z,y)=-(x,y).
\]
The differential in degree three on a chain
$(x,z,w,y)$, where $z$ precedes $w$ along a
minimizing geodesic, is
\[
 \partial_3(x,z,w,y)
 =-(x,w,y)+(x,z,y).
\]
Every pair of internal points on the same minimizing
geodesic can be related using such boundaries, taking
a common later point if necessary. There are no
degree-three boundaries relating interior points of
distinct minimizing geodesics, because all chains
of total length $d$ lie on a single minimizing
geodesic. Hence
\[
 C_{2,d}/\operatorname{im}\partial_3
 \cong\Z[S_g(u,v)]
\]
and the induced differential is minus augmentation.
Taking its kernel proves the result.
\end{proof}

For a minimizing-label set $S$ write
$I(S)=\ker(\Z[S]\to\Z)$. The proposition identifies
the local coefficient $I(S)$ with magnitude homology;
the inclusion $I(S)\hookrightarrow I$ is natural
after lifts have been chosen. This is the local-to-global
map used in the localization theorem below.

\begin{remark}
The endpoint-fixed statement is made only for $x\ne y$.
The diagonal has length zero and a different chain complex.
No identification of magnitude homology with the entire
simplicial nerve is asserted: the correspondence here
is specifically in degree two and at the minimal length.
\end{remark}
